% from ega0-21-fr.tex
\oldpage[0\textsubscript{IV-1}]{174}

\begin{corollary}[21.8.4]
\label{0.21.8.4-fr}
Soient $K$ un corps de caractéristique $p>0$, $k_0$ un sous-corps de $K$,
$(k_\alpha)_{\alpha\in I}$ une famille filtrante décroissante de sous-corps
de $K$, contenant $k_0$, telle que
\[
  \bigcap_{\alpha\in I}k_\alpha(K^p)=k_0(K^p).
\]
Soient $L$ une extension de $K$ telle que $\Upsilon_{L/K/k_0}$ soit un
$L$-espace vectoriel de rang fini~; alors, pour tout corps $K'$ tel que
$K\subset K'\subset L$, il existe $\alpha\in I$ tel que $k_\alpha$ soit
$k_0$-admissible pour l'extension $L$ de $K'$.
\end{corollary}

\begin{proof}
Il suffit d'appliquer (21.8.3) et (21.6.6).
\end{proof}

\begin{corollary}[21.8.5]
\label{0.21.8.5-fr}
Soient $K$ un corps de caractéristique $p>0$, $k_0$ un sous-corps de $K$,
$(k_\alpha)_{\alpha\in I}$ une famille filtrante décroissante de sous-corps
de $K$ contenant $k_0$, telle que
\[
  \bigcap_{\alpha\in I}k_\alpha(K^p)=k_0(K^p).
\]
Alors, pour toute extension $L$ de $K$ telle que $\Upsilon_{L/K/k_0}$ soit
un $L$-espace vectoriel de rang fini, on a
\[
  \bigcap_{\alpha\in I}k_\alpha(L^p)=k_0(L^p).
\]
\end{corollary}

\begin{proof}
Supposons en effet que $M$ soit une extension de $L$ telle que
$\Upsilon_{M/L/k_0}$ soit un $M$-espace vectoriel de rang fini. La suite
exacte (21.6.1.1)
\[
  0\longrightarrow \Upsilon_{L/K/k_0}\otimes_L M
  \longrightarrow \Upsilon_{M/K/k_0}
  \longrightarrow \Upsilon_{M/L/k_0}
\]
et l'hypothèse montrent alors que $\Upsilon_{M/K/k_0}$ est aussi un
$M$-espace vectoriel de rang fini. Il existe donc un indice $\alpha$ tel que
$k_\alpha$ soit $k_0$-admissible pour l'extension $M$ de $K$ (21.8.3), donc
aussi pour l'extension $M$ de $L$ (21.6.6)~; cela ayant lieu pour toute
extension $M$ de $L$ telle que $\Upsilon_{M/L/k_0}$ soit de rang fini, le
corollaire résulte de l'équivalence de $a)$ et $b)$ dans (21.8.3).
\end{proof}

\begin{corollary}[21.8.6]
\label{0.21.8.6-fr}
Soient $K$ un corps, $p$ son exposant caractéristique, $k_0$ un sous-corps de
$K$. Si $L$ est une extension de $K$ telle que $\Upsilon_{L/K/k_0}$ soit un
$L$-espace vectoriel de rang fini, il existe un sous-corps $k$ de $K$,
contenant $k_0(K^p)$, tel que $[K:k]$ soit fini, et qui soit
$k_0$-admissible pour l'extension $L$ de $K$.
\end{corollary}

\begin{proof}
Il suffit en effet, en vertu de (21.8.3), de construire une famille filtrante
décroissante $(k_\alpha)_{\alpha\in I}$ de sous-corps de $K$, contenant
$K^p$ et $k_0$, pour lesquels $[K:k_\alpha]<+\infty$ et
$\bigcap_\alpha k_\alpha=k_0(K^p)$. Pour cela on considère une $p$-base
$(x_\lambda)_{\lambda\in J}$ de $K$ sur $k_0$ et, pour toute partie finie
$H$ de $J$, on considère le sous-corps $k_H$ de $K$ engendré par $k_0(K^p)$
et les $x_\lambda$ d'indice $\lambda\in J-H$~; il résulte de cette définition
que $(x_\lambda)_{\lambda\in H}$ est une $p$-base de $K$ sur $k_H$, et on en
conclut aussitôt que les $k_H$ vérifient les conditions voulues.
\end{proof}

\begin{remarks}[21.8.7]
\label{0.21.8.7-fr}
\begin{enumerate}
  \item[(i)] On a déjà vu (21.7.2) que si $L$ est une extension de
  \emph{type fini} de $K$, $\Upsilon_{L/K/k_0}$ est de rang fini pour tout
  sous-corps $k_0$ de $K$. Il en est de même si $L$ est une extension
  \emph{séparable} de $K$, car en vertu de (20.6.19), on a
  $\Upsilon_{L/K/k_0}=0$. Enfin, si $L$ est une \emph{extension de type fini
  d'une extension séparable $L_0$ de $K$}, le même raisonnement que dans
  (21.8.5) montre que $\Upsilon_{L/K/k_0}$ est encore de rang fini (et en fait
  est isomorphe à un sous-espace de $\Upsilon_{L/L_0/k_0}$).

  